problem:  
Let \((M,\Delta,g_\Delta)\) be a complete, simply connected, equiregular sub‑Riemannian manifold with Hausdorff dimension \(Q\), and assume that for every \(p\in M\) the nilpotent approximation \(\mathrm{Nil}_p(M)\) is isomorphic to a fixed Carnot group \(G\).  
Define the \(k\)-dimensional filling function using horizontal integral currents by  
\[
\mathrm{Fill}_M^{(k)}(r)=
\sup_{\substack{\partial S=0,\\\mu_k(S)\le r}}
\inf_{\partial T=S}\mu_{k+1}(T),\qquad
\beta(k,M)=\limsup_{r\to\infty}\frac{\log\mathrm{Fill}_M^{(k)}(r)}{\log r}.
\]

It is known that for the model Carnot group \(G\), \(\mathrm{Fill}_G^{(k)}(r)\simeq r^{\beta(k,G)}\) for some exponent \(\beta(k,G)\) depending on \(k\) and the stratified Lie algebra of \(G\).

**Problem (Carnot‑rigidity of filling exponents):**  
Is it true that
\[
\beta(k,M)=\beta(k,G)\ \text{for all } k,
\]
whenever \(M\) is equiregular and all its nilpotentizations are isomorphic to \(G\)?  
In other words, does the local Carnot model fully determine the asymptotic filling behavior of \(M\)?  

If not, construct or characterize manifolds \(M\) with \(\mathrm{Nil}_p(M)\cong G\) everywhere but \(\beta(k,M)\neq\beta(k,G)\).  In particular, could analytic phenomena that are invisible to the nilpotent approximation—such as the structure or abundance of **abnormal geodesics**—alter the large‑scale filling exponent?

---

Why is it a "good" problem:  
- **Profound insight:**  
  The question targets the heart of sub‑Riemannian geometry: whether large‑scale geometric‐measure invariants are entirely determined by infinitesimal Carnot data.  A positive answer would establish a new rigidity principle analogous to curvature rigidity in Riemannian geometry; a negative answer would reveal that infinitesimal information is insufficient for predicting global asymptotic geometry.  

- **Bridge between fields:**  
  The problem connects **nilpotent approximation and homogeneity theory** (via Carnot models), **geometric measure theory** (through fillings and integral currents), and **control theory** (through abnormal geodesics and nonholonomy).  Any progress would hinge on unifying local analytic control structure with global geometric growth, bridging distinct areas of modern analysis and geometry.  

- **Extreme simplicity:**  
  The statement can be summarized in one line:  
  > “Do identical Carnot tangents force identical filling exponents?”  
  Yet its resolution would yield decisive information about how local nonholonomic degeneracies shape macroscopic quantitative geometry, illuminating a foundational aspect of sub‑Riemannian asymptotic behavior.